% Generated draft: requires mathematical and source review.
\documentclass[11pt]{article}
\usepackage[margin=25mm]{geometry}
\usepackage{fontspec}
\setmainfont{Latin Modern Roman}
\usepackage{amsmath,amssymb,mathtools,unicode-math}
\setmathfont{Latin Modern Math}
\usepackage{enumitem,needspace,xurl,hyperref}
\setlength{\parindent}{0pt}
\setlength{\parskip}{5pt}
\setlength{\emergencystretch}{3em}
\newcommand{\N}{\mathbb{N}}
\newcommand{\Z}{\mathbb{Z}}
\newcommand{\Q}{\mathbb{Q}}
\newcommand{\R}{\mathbb{R}}
\newcommand{\C}{\mathbb{C}}
\newcommand{\F}{\mathbb{F}}
\newcommand{\A}{\mathbb{A}}
\newcommand{\PP}{\mathbb{P}}
\newcommand{\E}{\mathbb{E}}
\newcommand{\eps}{\varepsilon}
\newcommand{\dd}{\,\mathrm{d}}
\begin{document}
\section*{UnsolvedMath ID: 30001538}
\textbf{Source problem code:} catalogue entry
\subsection{Definitions and mathematical statement}
The supplied record uses \(\rho_p\) for a \(p\)-adic representation and \(L_p(T)\) for an algebraic \(p\)-adic \(L\)-function, with \(p\) a prime. It does not identify the arithmetic objects to which these symbols refer, the domain and target of \(\rho_p\), the coefficient ring of \(L_p(T)\), or the meaning of its parameter \(T\). These symbols are therefore retained as unspecified objects from the source question.

\textbf{Research question.} Can analogues of the \(p\)-adic representation \(\rho_p\) and the algebraic \(p\)-adic \(L\)-function \(L_p(T)\) be constructed in knot theory?

The supplied record does not specify the class of knots, the ambient manifolds, the input data for either construction, or the properties that would qualify either object as an analogue. It also does not specify whether the constructions are sought for every prime, for a fixed prime, or under additional hypotheses. Consequently, this is a research direction whose precise existence statement remains unresolved at the formulation level; no particular representation, coefficient ring, knot invariant, or comparison property is imposed here.
\subsection{Short English statement}
The question asks for knot-theoretic counterparts of an unspecified arithmetic \(p\)-adic representation and algebraic \(p\)-adic \(L\)-function. The supplied record leaves the intended objects and analogy criteria undefined, so further source clarification is needed for a precise mathematical existence problem.
\subsection{Sources}
\begin{itemize}
\item Local registry: \texttt{data/unsolvedmath/problems.json}, permanent UnsolvedMath ID \(30001538\), source problem code \texttt{catalogue entry}, in the supplied Oberwolfach Reports Open Problems collection. No individual proposer is identified; the supplied proposed year is 2010.
\item Supplied report citation: \emph{Low-Dimensional Topology and Number Theory} (2011), \emph{Oberwolfach Reports}, pp.~2101--2163, DOI \texttt{10.4171/owr/2010/35}. The source URL and the literature-source link labelled ``original OWR report'' are \url{https://doi.org/10.4171/owr/2010/35}. The background describes these reports as documenting mathematical workshops and their open-problem discussions.
\item Supplied literature source labelled ``L2-signature theory'': \url{https://doi.org/10.4007/annals.2003.157.433}. No further bibliographic details are supplied for this reference.
\item Supplied literature assessment, checked at the exact date \texttt{2026-08-21}: the registry labels the problem open and describes existing \(L^2\)-signature and concordance theories as partial analogies, with no general construction answering the extracted question. This is supplied dated curation derived from the OWR record, not an independently verified literature assessment. Time-sensitive status and cited claims require independent verification.
\item The registry labels the statement \texttt{corrected\_verified}, explaining that the complete statement was preserved while mathematical notation was placed in delimiters. This supplied editorial label does not establish that the missing definitions or intended mathematical target have been verified.
\end{itemize}
\end{document}
